\chapter{Rational numbers, ratios, and fields}\label{ch:rationals}
\begin{readergoal}Construct fractions and learn why units, denominators, and field axioms matter in quantitative science.\end{readergoal}
\section{Fractions as equivalence classes}
Take integer pairs $(a,b)$ with $b\ne0$. Say $(a,b)\sim(c,d)$ when $ad=bc$. The resulting equivalence class represents the fraction $a/b$. Ratios like $1/2=2/4$ are different representations of the same rational number. Addition and multiplication take their familiar forms
\[\frac ab+\frac cd=\frac{ad+bc}{bd},\qquad \frac ab\cdot\frac cd=\frac{ac}{bd}.\]
Rationals form the field $\mathbb Q$: nonzero values possess multiplicative inverses, and the field operations obey associativity, commutativity, and distributivity. Division by zero is undefined: a rate with zero exposure duration is not an extreme number but an invalid operation.

\section{What does a ratio mean physically?}
Suppose a unit emits 75 spikes over 15 seconds. Its time-averaged rate is $75/15=5$ spikes/s. The numerator and denominator have different roles and units. When a software routine returns a mean across observations, ask what each observation represents and whether every observation has the same weight. The mean of within-trial rates need not equal total count divided by total exposure if trial durations differ.

\section{Rationals are dense but incomplete}
Between any two rational numbers $a<b$, their midpoint $(a+b)/2$ is another rational. Yet the equation $x^2=2$ has no rational solution. To show this, suppose $x=p/q$ is in lowest terms and $p^2=2q^2$. Then $p$ must be even; substituting $p=2k$ shows $q$ also even, contradicting coprimality. Thus the rational numbers have gaps from the viewpoint of limits and geometry.

\begin{workedexample}[Units and normalization]
Compare $\mathrm{spikes}/\mathrm{bin}$ with $\mathrm{spikes}/\mathrm{s}$. If the bin width is $\Delta t$ seconds, the rate estimate is $r=n/\Delta t$ in spikes/s. Changing $\Delta t$ changes the variance and resolution of the estimate even if the underlying experiment is unchanged. This anticipates the bias-variance trade-off in spike-rate estimation.
\end{workedexample}
\begin{exercises}
\item Explain what fails if $b=0$ is allowed in a fraction.
\item Calculate the pooled firing rate across a 5-s trial with 20 spikes and a 15-s trial with 30 spikes; compare with the unweighted average of the two trial rates.
\item Prove that between two unequal rationals there is a rational number.
\end{exercises}